% !TEX root = ../Presentation.tex
% ---------------------------------------------------------------------------
% Summary, open questions and next steps
% Slides taken from ../Presentation_ABC/sections; trimmed ones are marked.
% ---------------------------------------------------------------------------
\section{Summary}

% trimmed: elpd and p_loo rows replaced by a ranking
\begin{frame}{Comparison of the models}
  \centering
  \begin{tabular}{lccc}
    \toprule
     & A & B & C \\
    \midrule
    Sessions & pooled & hierarchical & hierarchical \\
    Distractor & separate groups & separate groups & covariate $\beta_2 x_\text{dist}$ \\
    Prediction (LOO) & 3rd & \textbf{best} & 2nd \\
    PSE shift, left / right ($\mu m$) & $-4.2$ / $-6.1$ & $-4.3$ / $-6.1$ & $-4.5$ / $-5.7$ \\
    \bottomrule
  \end{tabular}

  \vfill
  \raggedright
  \begin{itemize}
    \item All three models agree on the PSE shift
    \item B predicts best: sessions differ (B $\gg$ A), and lapse rates differ between
          conditions (B $>$ C)
  \end{itemize}
\end{frame}

\begin{frame}{Conclusions}
  \begin{itemize}
    \item The distractor lowers the PSE: by about 4.2--4.5~$\mu m$ (left) and 5.7--6.1~$\mu m$
          (right), in all three models.
    \item It raises the rate of ``high'' responses to low stimulus amplitudes on bimanual trials
          ($\gamma_h$); whether this grows with distractor amplitude is open.
    \item The sessions differ; a hierarchical model is needed.
    \item The shift grows with distractor amplitude, and a weak distractor may shift the PSE
          more than a proportional model allows (right hand).
  \end{itemize}
\end{frame}

% new for this audience
\begin{frame}{Open questions}
  \begin{itemize}
    \item \textbf{The extra ``high'' responses to low stimulus amplitudes on bimanual trials:} they seem
          to cluster in sessions with a strong distractor (being checked). What could produce
          them: responses to the distractor hand, attention, motor errors?
    \item \textbf{Amplitude and time:} the distractor amplitude rose over the sessions, so a
          growing effect could also be training or adaptation. Was the schedule chosen for a
          reason? Did anything else change over those months?
    \item \textbf{Left vs right:} the right hand shifts more, even at weak distractors. Is a
          hand asymmetry expected?
  \end{itemize}
\end{frame}

\begin{frame}{Next steps}
  \begin{itemize}
    \item A new Model D, revising Model C: lapse rates per condition, testing whether the bimanual
          $\gamma_h$ depends on distractor amplitude (or session order).
    \item A distractor-presence term only if it improves prediction.
    \item Model recovery for B against D.
  \end{itemize}
\end{frame}
